
\documentclass[11pt,a4paper]{article}

%% ─── Packages ───────────────────────────────────────────────────────────────
\usepackage[margin=0.85in, top=0.9in, bottom=0.9in]{geometry}
\usepackage[T1]{fontenc}
\usepackage[utf8]{inputenc}
\usepackage{lmodern}
\usepackage{setspace}
\usepackage{hyperref}
\usepackage{xcolor}
\usepackage{enumitem}
\usepackage{listings}
\usepackage{titlesec}
\usepackage{graphicx}
\usepackage{mdframed}
\usepackage{booktabs}
\usepackage{array}
\usepackage{parskip}

%% ─── Colours ─────────────────────────────────────────────────────────────────
\definecolor{codebg}{RGB}{248,248,252}
\definecolor{codeframe}{RGB}{210,214,230}
\definecolor{codecomment}{RGB}{80,140,80}
\definecolor{codestring}{RGB}{180,40,40}
\definecolor{codekeyword}{RGB}{20,60,170}
\definecolor{headerblue}{RGB}{30,60,130}
\definecolor{accentblue}{RGB}{41,98,200}
\definecolor{lightgray}{RGB}{245,246,250}
\definecolor{screengray}{RGB}{200,200,210}

%% ─── Code style ──────────────────────────────────────────────────────────────
\lstdefinestyle{ts}{
  backgroundcolor=\color{codebg},
  basicstyle=\ttfamily\footnotesize,
  frame=single,
  rulecolor=\color{codeframe},
  breaklines=true,
  showstringspaces=false,
  keywordstyle=\color{codekeyword}\bfseries,
  commentstyle=\color{codecomment}\itshape,
  stringstyle=\color{codestring},
  columns=fullflexible,
  numbers=none,
  xleftmargin=6pt,
  xrightmargin=6pt
}
\lstdefinestyle{html}{
  backgroundcolor=\color{codebg},
  basicstyle=\ttfamily\footnotesize,
  frame=single,
  rulecolor=\color{codeframe},
  breaklines=true,
  showstringspaces=false,
  keywordstyle=\color{codekeyword}\bfseries,
  stringstyle=\color{codestring},
  columns=fullflexible,
  numbers=none,
  xleftmargin=6pt,
  xrightmargin=6pt
}

%% ─── Hyperlinks ──────────────────────────────────────────────────────────────
\hypersetup{colorlinks=true, linkcolor=accentblue, urlcolor=accentblue}

%% ─── Spacing / lists ─────────────────────────────────────────────────────────
\setstretch{1.13}
\setlist[itemize]{noitemsep, topsep=3pt, leftmargin=*}
\setlist[enumerate]{noitemsep, topsep=3pt, leftmargin=*}

%% ─── Section headings ────────────────────────────────────────────────────────
\titleformat{\section}
  {\large\bfseries\color{headerblue}}
  {\thesection.}{0.5em}{}[\titlerule]
\titleformat{\subsection}
  {\normalsize\bfseries\color{accentblue}}
  {}{0em}{}
\titleformat{\subsubsection}
  {\small\bfseries}
  {}{0em}{}

%% ─── Screenshot box ──────────────────────────────────────────────────────────
\newmdenv[
  backgroundcolor=lightgray,
  linecolor=screengray,
  linewidth=1pt,
  innerleftmargin=8pt,
  innerrightmargin=8pt,
  innertopmargin=8pt,
  innerbottommargin=8pt,
  skipabove=6pt,
  skipbelow=6pt
]{screenshotbox}

%% ─── Document ────────────────────────────────────────────────────────────────
\begin{document}

%% ── Title block ──────────────────────────────────────────────────────────────
\begin{center}
  {\LARGE\bfseries\color{headerblue} MEAN Stack Practical File}\\[0.5em]
  {\large\bfseries Authenticated Todo Management Application}\\[0.3em]
  {\normalsize MongoDB \textbullet{} Express \textbullet{} Angular 19 \textbullet{} Node.js}
\end{center}

\noindent\rule{\linewidth}{0.8pt}

%%═══════════════════════════════════════════════════════════════════════════════
\section*{Practical 8 — Building a Single-Page Application with Angular}
%%═══════════════════════════════════════════════════════════════════════════════

\subsection*{Aim}
To build an Angular Single-Page Application (SPA) that demonstrates navigation, modular
component architecture, HTML content binding, routing with parameters, reactive forms, and
server-submitted form data handling.

\subsection*{Required Theory}
An SPA loads a single HTML shell and then dynamically swaps content via Angular's router and
component system.  Key Angular concepts used in this practical:

\begin{itemize}
  \item \textbf{Router \& Navigation} — \texttt{RouterOutlet}, \texttt{RouterModule},
        \texttt{router.navigate([])}
  \item \textbf{Routing Parameters} — \texttt{:id} path segment read via \texttt{ActivatedRoute}
  \item \textbf{Modular Architecture} — standalone components, lazy loading, barrel imports
  \item \textbf{Data Binding} — interpolation \texttt{\{\{\}\}}, property binding \texttt{[prop]},
        event binding \texttt{(event)}, and two-way binding \texttt{[(ngModel)]}
  \item \textbf{Template-Driven Forms} — \texttt{FormsModule}, \texttt{ngModel},
        \texttt{ngSubmit}
  \item \textbf{HTTP \& Services} — \texttt{HttpClient}, \texttt{BehaviorSubject} reactive stream
\end{itemize}

%% ── 8.1 App-Level Routing ────────────────────────────────────────────────────
\subsection*{1.\ Route Definitions — \texttt{app.routes.ts}}

The application defines four routes.  The \textbf{dashboard} and \textbf{task/:id} routes are
protected by an auth guard and loaded lazily.

\begin{lstlisting}[style=ts, language={[Sharp]C}]
export const routes: Routes = [
  {
    path: 'login',
    loadComponent: () =>
      import('./components/auth/auth.component')
        .then(m => m.AuthComponent)
  },
  {
    path: 'dashboard',
    loadComponent: () =>
      import('./components/todo-list/todo-list.component')
        .then(m => m.TodoListComponent),
    canActivate: [authGuard]        // route protection
  },
  {
    // Routing parameter: task detail page via /task/:id
    path: 'task/:id',
    loadComponent: () =>
      import('./components/task-detail/task-detail.component')
        .then(m => m.TaskDetailComponent),
    canActivate: [authGuard]
  },
  { path: '', redirectTo: 'dashboard', pathMatch: 'full' },
  { path: '**',  redirectTo: 'dashboard' }
];
\end{lstlisting}

%% ── 8.2 Root Component  ──────────────────────────────────────────────────────
\subsection*{2.\ Root Component — \texttt{app.ts}}

The root component acts as the shell.  It validates the stored JWT on every startup and
navigates the user accordingly.

\begin{lstlisting}[style=ts, language={[Sharp]C}]
@Component({
  selector: 'app-root',
  standalone: true,
  imports: [CommonModule, RouterOutlet],
  templateUrl: './app.html'
})
export class App implements OnInit {
  isCheckingSession = true;

  constructor(private authService: AuthService,
              private router: Router,
              private cdr: ChangeDetectorRef) {
    // Hydrate user from localStorage immediately so components
    // know a session exists before the async token check finishes.
    this.authService.hydrateSessionFromStorage();
  }

  ngOnInit(): void {
    this.authService.restoreSession().subscribe(user => {
      this.isCheckingSession = false;
      this.cdr.detectChanges();        // prevent NG0100 in SSR
      if (user) {
        if (this.router.url === '/login' || this.router.url === '/')
          this.router.navigate(['/dashboard']);
      } else {
        this.router.navigate(['/login']);
      }
    });
  }
}
\end{lstlisting}

%% ── 8.3 Routing Parameters ───────────────────────────────────────────────────
\subsection*{3.\ Reading Route Parameters — \texttt{task-detail.component.ts}}

The \texttt{task/:id} route passes the task ObjectId in the URL.
\texttt{ActivatedRoute.paramMap} is used to read the parameter and fetch the matching task.

\begin{lstlisting}[style=ts, language={[Sharp]C}]
export class TaskDetailComponent implements OnInit {
  task: Task | null = null;
  taskId: string | null = null;

  constructor(private route: ActivatedRoute,
              private router: Router,
              private todoService: TodoService) {}

  ngOnInit(): void {
    this.route.paramMap.subscribe(params => {
      this.taskId = params.get('id');     // read :id from URL
      if (!this.taskId) {
        this.router.navigate(['/dashboard']);
        return;
      }
      this.todoService.getTask(this.taskId).pipe(
        timeout(8000),
        catchError(() => {
          this.errorMessage = 'Task not found.';
          return of(null);
        })
      ).subscribe(task => { this.task = task; });
    });
  }
}
\end{lstlisting}

Navigation to the detail page from the list uses:
\begin{lstlisting}[style=ts, language={[Sharp]C}]
viewTaskDetail(task: Task): void {
  this.router.navigate(['/task', task._id]);   // /task/abc123
}
\end{lstlisting}

%% ── 8.4 Template-Driven Form ─────────────────────────────────────────────────
\subsection*{4.\ Template-Driven Form — \texttt{auth.component.html}}

Two-way binding via \texttt{[(ngModel)]} and form submission via \texttt{(ngSubmit)}:

\begin{lstlisting}[style=html]
<form class="auth-form" (ngSubmit)="submit()">
  <!-- Conditional field shown only for registration -->
  <div *ngIf="mode === 'register'" class="form-group">
    <label for="name">Name</label>
    <input id="name" name="name" [(ngModel)]="form.name"
           type="text" placeholder="Your name" />
  </div>

  <div class="form-group">
    <label for="email">Email</label>
    <input id="email" name="email" [(ngModel)]="form.email"
           type="email" placeholder="you@example.com" />
  </div>

  <div class="form-group">
    <label for="password">Password</label>
    <input id="password" name="password"
           [(ngModel)]="form.password" type="password" />
  </div>

  <p *ngIf="errorMessage" class="error-message">{{ errorMessage }}</p>

  <button type="submit" [disabled]="isSubmitting">
    {{ isSubmitting ? 'Please wait...'
       : (mode === 'login' ? 'Login' : 'Create Account') }}
  </button>
</form>
\end{lstlisting}

%% ── 8.5 Filter Form + Data Binding ──────────────────────────────────────────
\subsection*{5.\ Content Binding \& Filter Form — \texttt{todo-list.component.html}}

The dashboard demonstrates multiple binding techniques simultaneously:

\begin{lstlisting}[style=html]
<!-- Interpolation binding: show live stats -->
<span>Total: {{ getTotalCount() }}</span>
<span>Pending: {{ getPendingCount() }}</span>

<!-- Two-way binding on search input -->
<input type="text" placeholder="Search tasks..."
       [(ngModel)]="searchQuery" (input)="onFilterChange()" />

<!-- Property + event binding on select -->
<select [(ngModel)]="filterPriority"
        (change)="onFilterChange()">
  <option value="all">All Priorities</option>
  <option value="high">High</option>
</select>

<!-- *ngFor to render task list -->
<app-todo-item *ngFor="let task of filteredTasks"
  [task]="task"
  (onToggle)="toggleTask($event)"
  (onEdit)="openEditForm($event)"
  (onDelete)="deleteTask($event)">
</app-todo-item>
\end{lstlisting}

%% ── 8.6 TodoService — Reactive Stream ───────────────────────────────────────
\subsection*{6.\ Reactive Data Service — \texttt{todo.service.ts}}

A \texttt{BehaviorSubject} acts as the in-memory task store; all components subscribe to it
for instant UI updates without a full reload.

\begin{lstlisting}[style=ts, language={[Sharp]C}]
@Injectable({ providedIn: 'root' })
export class TodoService {
  private tasksSubject = new BehaviorSubject<Task[]>([]);
  public tasks$ = this.tasksSubject.asObservable();

  loadTasks(sortBy = 'createdAt'): void {
    const params = new HttpParams().set('sortBy', sortBy);
    // Auth header attached automatically by the interceptor
    this.http.get<Task[]>(this.apiUrl, { params })
      .pipe(tap(tasks => this.tasksSubject.next(tasks)))
      .subscribe();
  }

  createTask(task: Task): Observable<Task> {
    return this.http.post<Task>(this.apiUrl, task).pipe(
      tap(created =>
        this.tasksSubject.next(
          this.sortTasks([created, ...this.tasksSubject.value])
        )
      )
    );
  }
}
\end{lstlisting}

\subsection*{Output / Screenshots}

\begin{screenshotbox}
\textbf{[Screenshot 1]} Angular application running at \texttt{localhost:4200}\\[4pt]
\textit{Insert screenshot here — Dashboard with tasks list, filters and stats bar.}\\[60pt]
\rule{\linewidth}{0.4pt}
\end{screenshotbox}

\begin{screenshotbox}
\textbf{[Screenshot 2]} Task detail page accessed via routing parameter \texttt{/task/:id}\\[4pt]
\textit{Insert screenshot here.}\\[45pt]
\rule{\linewidth}{0.4pt}
\end{screenshotbox}

\newpage

%%═══════════════════════════════════════════════════════════════════════════════
\section*{Practical 9 — Secure User Authentication System in Angular}
%%═══════════════════════════════════════════════════════════════════════════════

\subsection*{Aim}
To implement a complete JWT-based user authentication system consisting of a login/register
form, token storage in \texttt{localStorage}, an HTTP interceptor that auto-attaches the token,
protected routes, and a logout function.

\subsection*{Required Theory}
\begin{itemize}
  \item \textbf{JWT (JSON Web Token)} — a self-contained signed token encoding user identity.
  \item \textbf{Bearer Authentication} — token sent in the \texttt{Authorization} HTTP header.
  \item \textbf{Angular HTTP Interceptor} — middleware that transforms every outgoing request.
  \item \textbf{Route Guard} — \texttt{CanActivateFn} that blocks unauthenticated navigation.
  \item \textbf{BehaviorSubject} — reactive store holding the logged-in user state.
\end{itemize}

\subsection*{Authentication Flow}

\begin{enumerate}
  \item User submits login credentials via Angular form.
  \item Frontend sends \texttt{POST /api/auth/login} to Express backend.
  \item Backend verifies password hash and returns \texttt{\{token, user\}}.
  \item \texttt{AuthService.storeSession()} saves token + user to \texttt{localStorage}.
  \item Angular interceptor reads the token and attaches it to every \texttt{/api/*} request.
  \item Auth guard checks \texttt{getToken()} before allowing access to \texttt{/dashboard}.
  \item Logout clears storage and redirects to \texttt{/login}.
\end{enumerate}

%% ── 9.1 Backend: login endpoint ─────────────────────────────────────────────
\subsection*{1.\ Node.js/Express — \texttt{backend/routes/auth.js}}

\begin{lstlisting}[style=ts, language={[Sharp]C}]
// POST /api/auth/login
router.post('/login', async (req, res) => {
  const { email, password } = req.body;
  const user = await User.findOne(
    { email: email.trim().toLowerCase() }
  );

  if (!user || !verifyPassword(password, user.passwordHash)) {
    return res.status(401)
              .json({ message: 'Invalid email or password' });
  }

  return res.json({
    token: generateToken(user),   // signs a JWT
    user:  sanitizeUser(user)     // strips passwordHash
  });
});

// POST /api/auth/register
router.post('/register', async (req, res) => {
  const { name, email, password } = req.body;
  const user = await User.create({
    name:         name.trim(),
    email:        email.trim().toLowerCase(),
    passwordHash: hashPassword(password)
  });
  return res.status(201).json({
    token: generateToken(user),
    user:  sanitizeUser(user)
  });
});
\end{lstlisting}

%% ── 9.2 Backend: auth middleware ─────────────────────────────────────────────
\subsection*{2.\ JWT Middleware — \texttt{backend/middleware/auth.js}}

All task routes pass through this middleware before executing:

\begin{lstlisting}[style=ts, language={[Sharp]C}]
export default async function authMiddleware(req, res, next) {
  const authHeader = req.headers.authorization;
  if (!authHeader?.startsWith('Bearer '))
    return res.status(401).json({ message: 'No token provided' });

  const token   = authHeader.slice(7);
  const payload = verifyToken(token);           // throws on invalid
  const user    = await User.findById(payload.sub)
                            .select('-passwordHash');
  if (!user)
    return res.status(401).json({ message: 'User not found' });

  req.user = user;   // attach user to request
  next();
}
\end{lstlisting}

%% ── 9.3 AuthService ──────────────────────────────────────────────────────────
\subsection*{3.\ Angular AuthService — \texttt{services/auth.service.ts}}

\begin{lstlisting}[style=ts, language={[Sharp]C}]
@Injectable({ providedIn: 'root' })
export class AuthService {
  private readonly tokenKey = 'todo-auth-token';
  private readonly userKey  = 'todo-auth-user';
  private userSubject = new BehaviorSubject<AuthUser | null>(null);
  readonly user$ = this.userSubject.asObservable();

  // ── Login: POST credentials, store response ──
  login(payload: { email: string; password: string })
                                     : Observable<AuthUser> {
    return this.http
      .post<AuthResponse>(`${this.apiUrl}/login`, payload)
      .pipe(
        tap(res  => this.storeSession(res)),   // save token+user
        map(res  => res.user)
      );
  }

  // ── Persist token + user to localStorage ──
  private storeSession(res: AuthResponse): void {
    localStorage.setItem(this.tokenKey, res.token);
    localStorage.setItem(this.userKey, JSON.stringify(res.user));
    this.userSubject.next(res.user);
  }

  // ── Read token safely (SSR-safe) ──
  getToken(): string | null {
    if (typeof window === 'undefined') return null;
    return localStorage.getItem(this.tokenKey);
  }

  // ── Logout: clear everything ──
  logout(): void {
    localStorage.removeItem(this.tokenKey);
    localStorage.removeItem(this.userKey);
    this.userSubject.next(null);
  }
}
\end{lstlisting}

%% ── 9.4 HTTP Interceptor ─────────────────────────────────────────────────────
\subsection*{4.\ HTTP Interceptor — \texttt{interceptors/auth.interceptor.ts}}

The interceptor clones every outgoing request targeting \texttt{/api/} and adds the
\texttt{Authorization} header automatically:

\begin{lstlisting}[style=ts, language={[Sharp]C}]
export const authInterceptor: HttpInterceptorFn = (req, next) => {
  const platformId = inject(PLATFORM_ID);

  // Skip on server — localStorage unavailable during SSR
  if (!isPlatformBrowser(platformId)) return next(req);

  const token = inject(AuthService).getToken();

  if (token && req.url.includes('/api/')) {
    const cloned = req.clone({
      setHeaders: { Authorization: `Bearer ${token}` }
    });
    return next(cloned);     // forward the modified request
  }

  return next(req);          // pass through unchanged
};
\end{lstlisting}

The interceptor is registered globally in \texttt{app.config.ts}:
\begin{lstlisting}[style=ts, language={[Sharp]C}]
export const appConfig: ApplicationConfig = {
  providers: [
    provideHttpClient(
      withFetch(),
      withInterceptors([authInterceptor])  // register globally
    ),
    provideRouter(routes),
    provideClientHydration(withEventReplay())
  ]
};
\end{lstlisting}

%% ── 9.5 Route Guard ──────────────────────────────────────────────────────────
\subsection*{5.\ Auth Guard — \texttt{guards/auth.guard.ts}}

\begin{lstlisting}[style=ts, language={[Sharp]C}]
export const authGuard: CanActivateFn = () => {
  const auth   = inject(AuthService);
  const router = inject(Router);

  if (auth.getToken()) return true;    // allow navigation

  // Redirect unauthenticated users to the login page
  router.navigate(['/login']);
  return false;
};
\end{lstlisting}

The guard is applied in the route table via \texttt{canActivate: [authGuard]} on
\texttt{/dashboard} and \texttt{/task/:id}.

%% ── 9.6 Logout ───────────────────────────────────────────────────────────────
\subsection*{6.\ Logout — \texttt{todo-list.component.ts}}

\begin{lstlisting}[style=ts, language={[Sharp]C}]
requestLogout(): void {
  this.authService.logout();      // clear token + user state
  this.todoService.clearTasks();  // wipe in-memory task list
  this.router.navigate(['/login']);
}
\end{lstlisting}

%% ── 9.7 Session Restore on Reload ────────────────────────────────────────────
\subsection*{7.\ Session Restore on Page Reload — \texttt{app.ts}}

\begin{lstlisting}[style=ts, language={[Sharp]C}]
ngOnInit(): void {
  // 1. Populate user$ immediately from localStorage
  this.authService.hydrateSessionFromStorage();

  // 2. Validate token against backend (async)
  this.authService.restoreSession().subscribe(user => {
    this.isCheckingSession = false;
    if (user) {
      // Token valid — stay on current page
    } else {
      // Token expired or invalid — force re-login
      this.todoService.clearTasks();
      this.router.navigate(['/login']);
    }
  });
}
\end{lstlisting}

\subsection*{Output / Screenshots}

\begin{screenshotbox}
\textbf{[Screenshot 3]} Login page — Angular form with email/password fields\\[4pt]
\textit{Insert screenshot here.}\\[60pt]
\rule{\linewidth}{0.4pt}
\end{screenshotbox}

\begin{screenshotbox}
\textbf{[Screenshot 4]} Successful login — redirected to protected \texttt{/dashboard}\\[4pt]
\textit{Insert screenshot here.}\\[45pt]
\rule{\linewidth}{0.4pt}
\end{screenshotbox}

\begin{screenshotbox}
\textbf{[Screenshot 5]} Browser DevTools — \texttt{localStorage} storing JWT token\\[4pt]
\textit{Insert screenshot here (Application → Local Storage).}\\[45pt]
\rule{\linewidth}{0.4pt}
\end{screenshotbox}

\begin{screenshotbox}
\textbf{[Screenshot 6]} Network tab — \texttt{Authorization: Bearer ...} header on API request\\[4pt]
\textit{Insert screenshot here (Network → /api/tasks → Request Headers).}\\[45pt]
\rule{\linewidth}{0.4pt}
\end{screenshotbox}

\begin{screenshotbox}
\textbf{[Screenshot 7]} Direct navigation to \texttt{/dashboard} while logged out — redirected to \texttt{/login}\\[4pt]
\textit{Insert screenshot here showing route guard in action.}\\[45pt]
\rule{\linewidth}{0.4pt}
\end{screenshotbox}

%%═══════════════════════════════════════════════════════════════════════════════
\section*{Result}
%%═══════════════════════════════════════════════════════════════════════════════

\textbf{Practical 8} — A fully functional Angular SPA was built demonstrating modular
standalone components, lazy-loaded routing, route parameters (\texttt{/task/:id}),
template-driven forms with two-way binding (\texttt{[(ngModel)]}), and a reactive data
service using \texttt{BehaviorSubject}.

\textbf{Practical 9} — A complete JWT authentication system was implemented.  The login
form posts credentials to the Express backend, the returned token is stored in
\texttt{localStorage}, an Angular HTTP interceptor automatically attaches it to every
\texttt{/api/*} request, an auth guard blocks unauthenticated users from protected routes,
and a logout function clears all session data and redirects to the login page.

\end{document}
